\documentclass{article}
\usepackage[T1]{fontenc}
\usepackage{lmodern}
\usepackage{graphicx}
\usepackage{amsmath,amssymb}
\usepackage[a4paper,margin=2cm]{geometry}
\usepackage[absolute,overlay]{textpos}
\setlength{\TPHorizModule}{1mm}
\setlength{\TPVertModule}{1mm}
\usepackage[indonesian]{babel}
\usepackage{hyperref}
\setlength{\emergencystretch}{2em}
% unit: Halaman 9

\begin{document}

% === Halaman 9 (python/native) ===

Contoh 2: Misalkan bobot-bobot yang membentuk barisan superincreasing 

adalah \{2, 3, 6, 13, 27, 52\}, dan diketahui bobot knapsack (M) = 70. Kita akan 

mencari b1, b2, …, b6 sedemikian sehingga


70= 2b1 + 3b2 + 6b3 + 13b4 + 27b5 + 52b6

 Caranya sebagai berikut:

1) Bandingkan 70 dengan bobot terbesar, yaitu 52. Karena 52  70, 

maka 52 dimasukkan ke dalam knapsack.  → b6 = 1

2) Bobot total sekarang menjadi 70 – 52 = 18. Bandingkan 18 dengan 

bobot terbesar kedua, yaitu 27. Karena 27 > 18, maka 27 tidak 

dimasukkan ke dalam knapsack. → b5 = 0

3) Bandingkan 18 dengan bobot terbesar berikutnya, yaitu 13. Karena 

13  18, maka 13 dimasukkan ke dalam knapsack. → b4 = 1

9

\end{document}
